\documentclass{article}
\usepackage[UTF8]{ctex}
\usepackage{amsmath,amssymb,amsfonts}
\usepackage{enumitem}
\usepackage{geometry}
\geometry{a4paper, margin=1in}
% === 错题本排版优化（不参与题目标识） ===
\setlength{\parindent}{0pt}
\setlength{\parskip}{0.35em plus 0.10em minus 0.08em}
\setlength{\emergencystretch}{3em}
\setlist{itemsep=0.25em, topsep=0.40em, parsep=0pt}
\allowdisplaybreaks[2]
% === 错题本排版优化结束 ===
\begin{document}
\section*{电子科技大学2026年高等代数真题}
\section*{160．电子科技大学2026年高等代数试卷}
一．填空题．每题5分，共30分．\\

\subsection{题目1}
1．矩阵 \(A=\left(\begin{array}{ccc}3 & 2 & 4 \\ 1 & 1 & -1 \\ -3 & 6 & 1\end{array}\right), M_{i j}\) 表示 \(|A|\) 中 \((i, j)\) 元素的余子式，则 \(\sum_{i, j=1}^{3} M_{i j}=\)

\subsection{题目2}
2．\(f(\lambda)\) 是首项系数为 1 的整系数多项式，若 \(f(2+\sqrt{3})=0\) ，则 \(f(5)=\)

\subsection{题目3}
3．线性空间 \(V\) 的维数为 \(30, \mathcal{A}, \mathcal{B}, \mathcal{C}\) 为 \(V\) 上的线性变换，且 \(\mathcal{A} V, \mathcal{B} V, \mathcal{C} V\) 的维数均为 24 ，那么 \(\mathcal{A} \mathcal{B} \mathcal{C} V\)的最小维数是 \(\_\_\_\_\)\\
\(\_\_\_\_\)

\subsection{题目4}
4．矩阵 \(A\) 的零化多项式为 \(A^{4}+5 I=4 A^{2}\) ，且 \(\operatorname{det}(A)=8, A^{*}\) 为 \(A\) 的伴随矩阵，\(A^{*}-2 I\) 的最小多项式为 \(f(x)\) 为 \(\_\_\_\_\)

\subsection{题目5}
5．特征多项式为 \(\lambda^{3}(\lambda-1)^{4}\) 的复矩阵按相似分类，可分为 \(\_\_\_\_\)类．

\subsection{题目6}
6．设 \(V\) 为12维欧氏空间，子空间 \(U, W\) 的维数分别为 3,4 ，则 \(\operatorname{dim}\left(U^{\perp} \cap W^{\perp}\right)\) 的最小值为 \(\_\_\_\_\)．

二．计算题．每题15分，共60分．\\

\subsection{题目7}
7．在三维欧氏空间 \(\mathbb{R}^{3}\) 中，定义内积

\[
(\alpha, \beta)=x_{1} y_{1}+x_{2} y_{2}+x_{3} y_{3}, \alpha=\left(x_{1}, x_{2}, x_{3}\right), \beta=\left(y_{1}, y_{2}, y_{3}\right) .
\]

设 \(u, v, w \in \mathbb{R}^{3}\) ，其模长分别为 \(1,2,4\) ，它们两两之间的夹角为 \(\frac{\pi}{3}\) 。

\begin{itemize}
  \item[（1）] 直接写出 \(u, v, w\) 的 Gram 矩阵 \(G=\left(\begin{array}{ccc}(u, u) & (u, v) & (u, w) \\ (v, u) & (v, v) & (v, w) \\ (w, u) & (w, v) & (w, w)\end{array}\right)\) ．
  \item[（2）] 求 \(\operatorname{det}\left(I+u^{\mathrm{T}} u+v^{\mathrm{T}} v+w^{\mathrm{T}} w\right)\) 的值．
\end{itemize}

\subsection{题目8}
8．已知 \(\beta=(2,1,7)^{\mathrm{T}}\) 可表示为 \(\alpha_{1}=(1,-2,4 c)^{\mathrm{T}}, \alpha_{2}=(2, c, 10)^{\mathrm{T}}, \alpha_{3}=(3,2,0)^{\mathrm{T}}\) 的两种不同系数的线性组合。

\begin{itemize}
  \item[（1）] 求 \(c\) 的值．
  \item[（2）] 求方程组 \(\left\{\begin{array}{l}x_{1}+x_{2}+x_{3}=6 \\ c x_{1}+2 x_{2}-x_{3}=5 \\ 3 x_{1}+5 x_{2}+(12-c) x_{3}=7\end{array}\right.\) 的通解．
\end{itemize}

\subsection{题目9}
9．令 \(\mathcal{A}\) 是3维线性空间 \(V\) 上的线性变换，\(\alpha \in V\) 满足 \(V\) 中任意向量可由 \(\alpha, \mathcal{A} \alpha, \mathcal{A}^{2} \alpha\) 线性表出，且\(2 \alpha-5 \mathcal{A} \alpha+4 \mathcal{A}^{2} \alpha-\mathcal{A}^{3} \alpha=0\) ．\\
（1）求 \(\mathcal{A}\) 的全部特征值与特征向量．\\
（2）是否存在 \(V\) 的一组基，使得 \(\mathcal{A}\) 在这组基下的矩阵为对角矩阵？若存在，求出此基，若不存在，请说明理由。

\subsection{题目10}
10．设二次型 \(f\left(x_{1}, x_{2}, x_{3}\right)=2 x_{1}^{2}+x_{2}^{2}+x_{3}^{2}+2 x_{1} x_{2}+2 x_{1} x_{3}\) ．

\begin{itemize}
  \item[（1）] 求 \(f(x)\) 的矩阵的特征值与特征向量．
  \item[（2）] 求正交线性替换 \(X=C Y\) ，使得 \(f\) 化为标准形．
  \item[（3）] 求 \(f\left(x_{1}, x_{2}, x_{3}\right)\) 在条件 \(x_{1}^{2}+x_{2}^{2}+x_{3}^{2}=4\) 下的最大值．
\end{itemize}
\end{document}
